problem:  
Let \((M^n,g)\) be a compact Einstein manifold with \(\mathrm{Ric}_g = (n-1)g\), and let \(\Delta_p\) denote the Hodge–Laplacian acting on coexact \(p\)-forms.  
By the Gallot–Meyer inequality, one has for any such manifold  
\[
\lambda_1^{(p)}(M,g) \ge (p+1)(n-p),
\]
with equality achieved for the canonical round sphere \(S^n\).  

Suppose \((M^n,g)\) satisfies for some small \(\varepsilon > 0\),
\[
\lambda_1^{(p)}(M,g) \le (p+1)(n-p) + \varepsilon,
\]
and that the dimension of the eigenspace \(E_{\lambda_1^{(p)}}(M,g)\) differs from that of the corresponding eigenspace on \(S^n\) by at most \(C(n,p)\,\varepsilon\).  

**Problem:**  
Does such a *spectral pinching* condition imply that \((M^n,g)\) is *quantitatively close* to the round sphere?  

More precisely, is there a universal function \(\Psi(\varepsilon)\to 0\) as \(\varepsilon \to 0\) such that  
\[
d_{GH}\big((M,g),(S^n,g_{S^n})\big)\le \Psi(\varepsilon)\,?
\]  
Alternatively, can one construct a sequence of compact Einstein manifolds with \(\mathrm{Ric}=(n-1)g_i\) satisfying  
\[
\lambda_1^{(p)}(M_i,g_i)\to (p+1)(n-p), \quad \dim E_{\lambda_1^{(p)}}(M_i,g_i)\to \dim E_{\lambda_1^{(p)}}(S^n),
\]
for which \((M_i,g_i)\) fail to converge in the Gromov–Hausdorff sense to the sphere?  

In other words, does *Obata-type rigidity for \(p\)-forms* extend beyond the scalar case, or does new geometric flexibility appear when the Hodge–Laplacian is viewed on middle-degree forms?

---

Why is it a "good" problem:  
This refined problem builds upon the foundational Obata rigidity phenomenon but extends it to the realm of Hodge theory, explicitly anchored to the known Gallot–Meyer inequality — ensuring that the “spherical eigenvalue” serves as the correct comparator in the Einstein class.  

It is a *good problem* because:  
- It addresses a **conceptual frontier**: extending scalar spectral rigidity to form-level structure, linking eigenvalues of \(\Delta_p\) to global geometry.  
- It forges a **bridge between spectral geometry, Einstein manifolds, and representation theory**, including the delicate multiplicity data that encode symmetry.  
- It has **extreme simplicity in form**—a one-line spectral inequality pinched near equality—but profound implications touching rigidity, Bochner–Weitzenböck identities, and potential classification of “nearly symmetric” Einstein spaces.  
- It remains **entirely open**: no known theory establishes quantitative rigidity or counterexamples for \(p\)-form eigenvalue pinching, yet its verification would be a natural “second-generation Obata theorem” enriching the unity between curvature, spectrum, and topology.